\chapter*{Résumé}
\addcontentsline{toc}{chapter}{Résumé}

\begin{otherlanguage}{french}
Les nomenclatures logicielles soutiennent la transparence de la chaîne
d'approvisionnement logicielle et la gestion des vulnérabilités. Cependant,
la plupart des méthodes de génération reposent sur le code source, les
manifestes de construction ou les métadonnées d'un gestionnaire de paquets.
Dans un contexte d'audit Android, le seul artefact disponible peut être une
\APK compilée, dans laquelle les composants tiers sont fusionnés au bytecode
\DEX puis transformés par la réduction, l'optimisation et l'obfuscation.

Ce mémoire d'ingénieur présente une chaîne statique, fondée sur des preuves et
contrôlée par des seuils, qui transforme une \APK en \SBOM. Elle extrait les
preuves provenant de l'archive, du manifeste, du bytecode, des ressources et
des fichiers natifs; compare les caractéristiques de l'APK à des profils de
référence Maven versionnés; sépare l'identité d'un composant de la confiance
dans sa version exacte; conserve les alternatives ambiguës; produit un fichier
CycloneDX JSON muni d'identifiants Package URL; et interroge OSV en tenant
compte du niveau de confiance. Les arbres de dépendances du code source servent
à l'évaluation lorsqu'ils sont disponibles, tandis que l'apprentissage
automatique et l'assistance liée à R8 restent limités au rapport.

L'évaluation distingue les applications contrôlées, les APK réelles validées
par le code source et classées comme propres ou partielles, les contrôles sans
composant Maven émis, la couverture du dénominateur source complet, la
reconnaissance majorée avec versions source connues, les expériences rejetées
et les cas limites. Le système retrouve 84/84 versions attendues dans les
applications contrôlées. Pour les 102 APK réelles validées par le code source
qui ont émis des composants Maven, 927/1\,837 versions émises sont confirmées,
soit un taux de confirmation des versions émises de 50,5\%, et non un rappel
de toutes les dépendances. Quatre APK supplémentaires validées par le code
source n'émettent aucun composant Maven et restent des contrôles de rappel.
Par rapport aux 12\,082 lignes de dépendances résolues depuis le code source,
927 lignes exactes sont émises de bout en bout. Une voie distincte, contrainte
par les versions source, confirme 6\,772/6\,772 lignes émises et n'est
présentée que comme une borne supérieure.

Le MVP obtenu démontre la faisabilité pratique d'une reconstruction de SBOM à
partir d'un binaire, tout en exposant ses limites concernant la couverture de
la base de référence, le rappel du moteur de correspondance, les versions
voisines, les familles d'artefacts, le code natif, le code dissimulé, le
passage à l'échelle et l'autorité de l'IA. Son principe central consiste à
préférer une connaissance partielle, vérifiable et accompagnée d'une
abstention explicite à une certitude non justifiée sur la version exacte.
\end{otherlanguage}

\vspace{0.5cm}
\noindent\textbf{Mots-clés ---}
\begin{otherlanguage}{french}
sécurité Android, analyse d'APK, nomenclature logicielle, détection de
bibliothèques tierces, attribution de composants, CycloneDX, gestion des
vulnérabilités, R8.
\end{otherlanguage}
